\section*{Chapitre \ref{ch:g11:redox} --- Oxydation et réduction}

\begin{solution}{exo:g11:redox:1}
\omterm{def:g11:redox:oxidant}{Oxydant}~: une espèce capable de gagner des \omterm{def:g9:inside-the-atom:nucleus}{électrons}. \omterm{def:g11:redox:oxidant}{Réducteur}~: une
espèce capable de céder des \omterm{def:g9:inside-the-atom:nucleus}{électrons}. \omterm{def:g11:redox:oxidation}{Oxydation}~: une perte
d'\omterm{def:g9:inside-the-atom:nucleus}{électrons}. \omterm{def:g11:redox:oxidation}{Réduction}~: un gain d'\omterm{def:g9:inside-the-atom:nucleus}{électrons}.
\end{solution}

\begin{solution}{exo:g11:redox:2}
\ce{Fe^{2+}}/\ce{Fe}~; \ce{Ag+}/\ce{Ag}~; \ce{I2}/\ce{I-}~;
\ce{Fe^{3+}}/\ce{Fe^{2+}}.
\end{solution}

\begin{solution}{exo:g11:redox:3}
\ce{Ag+ + e- <=> Ag}~; \ce{Al^{3+} + 3e- <=> Al}~;
\ce{Fe^{3+} + e- <=> Fe^{2+}}~; \ce{Cl2 + 2e- <=> 2Cl-}.
\end{solution}

\begin{solution}{exo:g11:redox:4}
Une \omterm{def:g11:redox:oxidation}{oxydation} (des \omterm{def:g9:inside-the-atom:nucleus}{électrons} sont perdus)~; le fer joue le rôle de
\omterm{def:g11:redox:oxidant}{réducteur}.
\end{solution}

\begin{solution}{exo:g11:redox:5}
Le zinc est oxydé et c'est le \omterm{def:g11:redox:oxidant}{réducteur}~; \ce{Cu^{2+}} est réduit et
c'est l'\omterm{def:g11:redox:oxidant}{oxydant}. Deux \omterm{def:g9:inside-the-atom:nucleus}{électrons} par \omterm{def:g7:atoms-and-molecules:atom}{atome} de zinc.
\end{solution}

\begin{solution}{exo:g11:redox:6}
On part de \ce{H2O2 -> H2O}~: % ce-unbalanced-ok (step 1 of the method)
oxygène~: deux O à gauche, donc \ce{2H2O} à droite~; hydrogène~: 2 à
gauche, 4 à droite, donc \ce{2H+} à gauche~; charge~: $+2$ à gauche, 0 à
droite, donc deux \omterm{def:g9:inside-the-atom:nucleus}{électrons} à gauche~:
\ce{H2O2 + 2H+ + 2e- <=> 2H2O}.
\end{solution}

\begin{solution}{exo:g11:redox:7}
\ce{I2 + 2e- -> 2I-} et \ce{2S2O3^{2-} -> S4O6^{2-} + 2e-}~; deux
\omterm{def:g9:inside-the-atom:nucleus}{électrons} chacune, donc \ce{I2 + 2S2O3^{2-} -> 2I- + S4O6^{2-}}.
\end{solution}

\begin{solution}{exo:g11:redox:8}
\ce{Cu -> Cu^{2+} + 2e-} et $2\times$ (\ce{Ag+ + e- -> Ag})~:
\ce{Cu + 2Ag+ -> Cu^{2+} + 2Ag}. Les \omterm{def:g9:ions:ion}{ions} \ce{Cu^{2+}} formés sont
bleus.
\end{solution}

\begin{solution}{exo:g11:redox:9}
\ce{C6H8O6 -> C6H6O6 + 2H+ + 2e-} et \ce{I2 + 2e- -> 2I-}~:
\ce{C6H8O6 + I2 -> C6H6O6 + 2I- + 2H+}. La vitamine C est le \omterm{def:g11:redox:oxidant}{réducteur}.
\end{solution}

\begin{solution}{exo:g11:redox:10}
\ce{Al -> Al^{3+} + 3e-} cède 3 \omterm{def:g9:inside-the-atom:nucleus}{électrons}, \ce{Cu^{2+} + 2e- -> Cu} en
capte 2~; le plus petit multiple commun est 6, donc la première est
multipliée par 2 et la seconde par 3~: \ce{2Al + 3Cu^{2+} -> 2Al^{3+} + 3Cu}.
\end{solution}

\begin{solution}{exo:g11:redox:11}
\ce{H2O2 + 2H+ + 2e- -> 2H2O} et $2\times$ (\ce{Fe^{2+} -> Fe^{3+} + e-})~:
\ce{H2O2 + 2H+ + 2Fe^{2+} -> 2H2O + 2Fe^{3+}}.
\end{solution}

\begin{solution}{exo:g11:redox:12}
$n(\ce{Zn}) = 1.30/65.4 = \qty{1.99e-2}{mol}$. Un \omterm{def:g7:atoms-and-molecules:atom}{atome} de cuivre se
dépose par \omterm{def:g7:atoms-and-molecules:atom}{atome} de zinc oxydé, donc $n(\ce{Cu}) = \qty{1.99e-2}{mol}$
et $m(\ce{Cu}) = 1.99\times 10^{-2} \times 63.5 = \qty{1.26}{g}$.
\end{solution}

\begin{solution}{exo:g11:redox:13}
\ce{Cr2O7^{2-} + 14H+ + 6e- -> 2Cr^{3+} + 7H2O} et
$6\times$ (\ce{Fe^{2+} -> Fe^{3+} + e-})~:
\ce{Cr2O7^{2-} + 14H+ + 6Fe^{2+} -> 2Cr^{3+} + 6Fe^{3+} + 7H2O}.
Six \ce{Fe^{2+}} par \omterm{def:g9:ions:ion}{ion} dichromate~: $6 \times 0.010 = \qty{0.060}{mol}$.
\end{solution}

\begin{solution}{exo:g11:redox:14}
\ce{2ClO- + 4H+ + 2e- -> Cl2 + 2H2O} et \ce{2Cl- -> Cl2 + 2e-}~; en
additionnant et en divisant par 2~: \ce{ClO- + Cl- + 2H+ -> Cl2 + H2O}.
Un acide fournit les \omterm{def:g9:ions:ion}{ions} \ce{H+}, et la réaction libère du dichlore, un
gaz toxique~: d'où l'avertissement de ne jamais \omterm{def:g2:mixing-and-dissolving:mix}{mélanger} les deux
produits.
\end{solution}

\begin{solution}{exo:g11:redox:15}
Les \omterm{def:g9:inside-the-atom:nucleus}{électrons} ne peuvent pas circuler librement dans la solution
(\cref{prop:g11:redox:no-free-electrons})~: ceux que perdent les \omterm{def:g7:atoms-and-molecules:atom}{atomes}
de zinc ne peuvent être captés que par un \omterm{def:g9:ions:ion}{ion} \ce{Cu^{2+}} qui touche le
\omterm{def:g9:periodic-table-first-look:metal}{métal}. L'\omterm{def:g7:atoms-and-molecules:atom}{atome} de cuivre se forme donc à la surface de la lame, et y
reste.
\end{solution}

\begin{solution}{pb:g11:redox:1}
\textbf{1.} \ce{Cr2O7^{2-}}/\ce{Cr^{3+}} et \ce{CH3COOH}/\ce{CH3CH2OH}.

\textbf{2.} L'éthanol est oxydé~: c'est le \omterm{def:g11:redox:oxidant}{réducteur}.

\textbf{3.} \ce{Cr2O7^{2-} + 14H+ + 6e- <=> 2Cr^{3+} + 7H2O}.

\textbf{4.} Le carbone est ajusté (deux de chaque côté). Oxygène~: deux à
gauche, un à droite, donc \ce{H2O} à droite. Hydrogène~: quatre à
gauche, $6 + 2 = 8$ à droite, donc \ce{4H+} à gauche. Charge~: $+4$ à
gauche, 0 à droite, donc quatre \omterm{def:g9:inside-the-atom:nucleus}{électrons} à gauche~:
\ce{CH3COOH + 4H+ + 4e- <=> CH3CH2OH + H2O}.

\textbf{5.} Douze \omterm{def:g9:inside-the-atom:nucleus}{électrons}~: $2\times$ la première, $3\times$ la seconde
renversée~:
\ce{2Cr2O7^{2-} + 3CH3CH2OH + 16H+ -> 4Cr^{3+} + 3CH3COOH + 11H2O}.

\textbf{6.} Les \omterm{def:g9:ions:ion}{ions} \ce{Cr2O7^{2-}} orange sont réduits en \omterm{def:g9:ions:ion}{ions}
\ce{Cr^{3+}} verts partout où l'éthanol les atteint.

\textbf{7.} $M = 2\times 39.1 + 2\times 52.0 + 7\times 16.0 =
\qty{294.2}{g/mol}$.

\textbf{8.} $n = 2.0\times 10^{-3}/294.2 = \qty{6.8e-6}{mol}$.

\textbf{9.} Trois \omterm{def:g7:atoms-and-molecules:molecule}{molécules} d'éthanol pour deux \omterm{def:g9:ions:ion}{ions} dichromate~:
$n = \tfrac{3}{2}\times 6.8\times 10^{-6} = \qty{1.0e-5}{mol}$.

\textbf{10.} $M(\ce{C2H6O}) = \qty{46.0}{g/mol}$, donc
$m = 1.02\times 10^{-5}\times 46.0 = \qty{4.7e-4}{g}$, environ
\qty{0.47}{mg}.

\textbf{11.} $0.25/0.47 \approx 0.53$~: environ \qty{53}{\percent} du
dichromate.

\textbf{12.} $0.53 \times 10 \approx \qty{5.3}{cm}$.

\textbf{13.} $\qty{0.25}{mg} \times 2100 = \qty{525}{mg}$ par litre de
sang, environ \qty{0.53}{g/L}.

\textbf{14.} L'\omterm{def:g4:air-a-mixture-of-gases:air}{air} expiré rencontre d'abord les grains de l'entrée~; le
dichromate y est en excès et retient tout l'éthanol, si bien que la
réaction est complète dans une première portion du tube, qui vire
entièrement au vert, avant que l'éthanol n'atteigne la suite.

\textbf{15.} La réaction est totale~: les \omterm{def:g9:ions:ion}{ions} \ce{Cr^{3+}} verts ne
redeviennent pas du dichromate, si bien qu'un tube usagé ne peut plus
mesurer.

% ledger: ghs:K2Cr2O7 (H350 among its hazard statements)
\textbf{16.} Il est très toxique (GHS06) et présente un grave danger pour
la santé (GHS08~: parmi ses mentions de danger figure «~peut provoquer le
cancer~»), en plus d'être corrosif et toxique pour les organismes
aquatiques~; un appareil destiné au public ne doit pas en contenir.

\textbf{17.} \ce{O2 + 4H+ + 4e- -> 2H2O} et
\ce{CH3CH2OH + H2O -> CH3COOH + 4H+ + 4e-}~:
\ce{CH3CH2OH + O2 -> CH3COOH + H2O}.

\textbf{18.} Quatre \omterm{def:g9:inside-the-atom:nucleus}{électrons}.

\textbf{19.} $n = 0.25\times 10^{-3}/46.0 = \qty{5.4e-6}{mol}$
d'éthanol~; \omterm{def:g9:inside-the-atom:nucleus}{électrons}~: $4n = \qty{2.2e-5}{mol}$, soit
$2.2\times 10^{-5}\times 6.02\times 10^{23} = \num{1.3e19}$ \omterm{def:g9:inside-the-atom:nucleus}{électrons}.

\textbf{20.} $Q = 1.3\times 10^{19}\times 1.60\times 10^{-19} \approx
\qty{2.1}{C}$. Chaque \omterm{def:g7:atoms-and-molecules:molecule}{molécule} d'éthanol envoie exactement quatre
\omterm{def:g9:inside-the-atom:nucleus}{électrons} dans le fil~: la charge est donc proportionnelle à la quantité
d'éthanol, et la mesurer revient à mesurer l'éthanol.
\end{solution}
